% ============================================================
\section{Conclusion}
\label{sec:conclusion}
% ============================================================

% Three paragraphs as of 20 Sep 2026 -- result, why no cheap diagnostic replaces the intervention, what
% to do -- funded by the cuts in S4, S6 and S7 this round. It was one paragraph for six rounds because
% each paragraph break costs a printed line the page limit did not have; the comment recording that is
% deleted rather than kept, because the constraint no longer applies and a stale justification is worse
% than none.
We set out to measure how much fine-tuning damages a foundation model's
pre-trained capability, and on these benchmarks it is mostly not
\emph{measurably} there to damage: no cell beats every admissible rung of an
% The middle clause -- four of the seven cells that beat a fitted ridge are improved by full fine-tuning
% in every seed -- is cut from this recap on page 9, where the body spills. It is stated in the abstract,
% in S1's result paragraph, in contribution 1, in S5.1 cell by cell and in S5.3's clause walk; the two
% clauses kept here are the ones this paragraph's opening sentence is about.
eight-baseline ladder, and no cell shows decisive
evidence of harmful encoder adaptation: at these seed counts 2 of 31 favour
freezing at a multiplicity-adjusted $q{<}0.05$, and both are cells the screen
rejects.
% The break between this paragraph and the next is closed on 21 Sep 2026: each paragraph break in this
% section costs a printed line, as the comment above records, and \S\ref{sec:exp4:poscontrol} needs
% them. Two paragraphs -- the result and why no cheap diagnostic replaces the intervention, then the
% advice -- rather than three. The break before "The advice therefore inverts" is the one that carries
% structure for a skimming reader, so it is the one kept.
The reason no cheap diagnostic rescues this is structural rather than statistical.
CKA tells us that the representation changed; only the intervention (fine-tune,
then fine-tune again with the encoder frozen) tells us whether that change was
valuable, and \textbf{across the cells we evaluate the former does not reliably
predict the latter}.
% The LOCO sentence -- "there is no association to work with either: under an analysis fixed before it
% was run we find no evidence that CKA improves out-of-sample prediction once backbone identity,
% horizon and n are known" -- is cut on 2026-09-24 for the four printed lines page 10 needs, and it is
% the right one to cut because the ABSTRACT states the same result with strictly more information (the
% 15 folds and the +0.201 -> +0.176 drop, which this version omitted), S6 argues it, and app:loco
% carries the pre-registered analysis. This was a pointer to a finding stated better twice elsewhere,
% which is the only kind of whole sentence left on this page that is not the sole site of its claim.
The claim we would defend outside this setting is
that methodological one: \textbf{a post-treatment similarity statistic cannot be
assumed to identify the value of a treatment that changed the representation.}
% Review M7, conclusion half: one sentence, so that the paper's own reading of its contribution is a
% falsification of a shortcut in use rather than a claimed discovery. The clause that stood here --
% "the news is not that the substitution is unjustified but that it is made" -- is cut on 2026-09-25
% for page 10's last two lines. \S\ref{sec:intro} states it in strictly stronger form: it concedes the
% result is unsurprising, names WHO makes the substitution (the prescriptions it has just quoted, and
% an earlier draft of this work), and says what measuring the cost takes. This clause was the
% compressed version of that, so the label and the price survive here and the argument survives there.
That is a falsification and not a discovery, and here we can price it: the best
threshold misplaces 5 of the 31 cells.
% The clause cut from here on 21 Sep 2026 -- "drift is measured AFTER the treatment and on the object
% the treatment altered, so even a strong association would not license a decision rule" -- restated
% the bold sentence two lines below it, which says the same thing in the general form the paper wants
% defended, and S3's post-treatment paragraph beside Eq. 1 argues it at length. Cut for
% \S\ref{sec:exp4:poscontrol}'s lines; the clause was the one whole sentence in this section that
% added no claim its neighbours do not.

The advice therefore inverts the literature's.  Check that
the pre-trained features beat a baseline fitted on your own data before spending
effort on LoRA, EWC or L2-SP, whose representation preservation is no guide to
their utility (\S\ref{sec:forecasting:replication}); if they do, settle the question against a frozen
encoder scored on windows you did not select on, in that order, because
\textbf{a forgetting result measured on cells with no measurable pre-trained
advantage is a statement about the benchmark}.  We do \emph{not} claim that drift
is generally unpredictive across backbones, which three cannot establish
(Appendix~\ref{app:claims}).
% The two POSITIVE claims were enumerated here as well -- CKA-does-not-identify, and
% the-screen-locates-cells -- and are cut on 21 Sep 2026 for \S\ref{sec:exp4:poscontrol}'s last two
% lines. Both are stated in the sentences immediately above (the bold spine sentence and the bold
% benchmark-selection sentence) and tabulated with their evidence and strength in app:claims, which is
% the section a reviewer asking for the three claims to be separated was pointed at. The DISCLAIMER is
% what is kept, because nothing else in the body says which reading we decline.
% The closing line -- "the cheap rule we hoped to report rests on a premise these benchmarks do not
% supply" -- is cut for the one printed line S5's treatment paragraph needs. It restated the
% benchmark-selection lesson three sentences above it rather than adding a claim, and it was the only
% whole sentence on the spilling page that did.
